\scanpage{107}
\mspage{22}
% archive sheet number: 106
% cancelled opening, source lines 1--3
\deleted{4.3. \unsure{Soit} Soit $X$ un schéma régulier, et soit}

\deleted{Rappelons \unsure{pour mémoire} une forme du th. de Abhyankar}

% source lines 4--8
\unsure{4.4.|4.3.} Soit $X$ un schéma régulier, $(\unsure{D_i|Z_i})_{1\leq i\leq r}$ une famille de \deleted{\unsure{sché.} diviseurs premiers rég.} sous-schémas réguliers irréductibles dans $X$, tels que le diviseur $D=\sum D_i$ soit un diviseur à croisements normaux (SGA 5 I \ldots).

% source lines 9--13
Soit $U=X-\operatorname{supp}D=X-\bigcup D_i$, \inserted{$f:U\longrightarrow X$ l'inclusion,} $n$ un entier premier aux caractéristiques résiduelles de $X$ en les points de \deleted{\unsure{$Z$}} $\operatorname{supp}D$, $\Lambda=\unsure{\mathbb Z/n\mathbb Z|Z/nZ}$. On se propose de calculer les faisceaux $R^i f_*(\Lambda_U)$ sur $X$. \deleted{$\unsure{R^i}f_*(\Lambda_U)$}

% source lines 14--15
\uline{Proposition} 4.5. \deleted{Avec} Les notations \inserted{étant celles} de 4.4, \deleted{et \unsure{supposons} que $X$ satisfait au th. de pureté}

(i) \deleted{$R^1 f_*(\Lambda_U)$ est canoniquement isom.} \inserted{On a un isom. canonique}

\[
R^1f_*(\Lambda_U)\simeq\deleted{\unsure{\text{le}}}\coprod_{1\leq i\leq r}(\mu_n^{\unsure{\otimes(-1|\otimes(-1)}})_{D_i}\simeq\coprod_{1\leq i\leq r}\Lambda_{D_i}(-1)
\tag*{(4.\deleted{4}5.1)}
\]

% boxed left-margin insertion keyed to (ii)
(ii) \inserted{Supposons que $X$ satisfait au th. de pureté cohomologique (4.1)} \inserted{Pour tout entier $p$,} L'homomorphisme \inserted{suivant,} déduit des cup-produits est un \deleted{\unsure{produit}} isomorphisme

\[
\bigwedge^p R^1 f_*(\Lambda_U)\xrightarrow{\sim}R^p f_*(\Lambda_U).
\tag*{(4.\deleted{4}5.2)}
\]

(iii). Par suite, \inserted{sous l'hyp. de (ii),} on a un isomorphisme canonique

\[
\begin{aligned}
R^p f_*(\Lambda_U)&\simeq\coprod_{1\leq i_1<i_2<\cdots<i_p\leq r}(\deleted{\unsure{\mu_n^{\otimes(-p)}}}\Lambda_X(-p))_{D_{i_1\cdots i_p}},\\
&\simeq\{\deleted{\unsure{\mu_n^{\otimes(-p)}}}\Lambda_X(-p)\}\otimes\coprod_{1\leq i_1<\cdots<i_p\leq r}\Lambda_{D_{i_1\cdots i_p}},
\end{aligned}
\tag*{(4.5.3)}
\]

% cancelled trial coefficient above the right part of (4.5.3)
\margin{\deleted{$\unsure{(\mu_n^{\otimes p})_X}$}}
% lower-left margin note
\margin{écrire $\Lambda_X(-p$}
% source last line, after a cancelled beginning
\deleted{\unsure{si} \illegible{3} \unsure{posé}} où on a posé $D_{i_1\cdots i_p}=D_{i_1}\unsure{\cap|\wedge}D_{i_2}\cap\cdots\cap D_{i_p}$.
